\documentclass[10pt]{article}
\usepackage[margin=0.75in]{geometry}
\usepackage{xcolor}
\usepackage{tcolorbox}
\usepackage{hyperref}
\usepackage{enumitem}
\usepackage{booktabs}

% Define colors
\definecolor{osaorange}{RGB}{220,115,38}
\definecolor{correctcolor}{RGB}{0,150,0}

% Configure hyperref
\hypersetup{
    colorlinks=true,
    linkcolor=blue,
    urlcolor=blue
}

% Compact spacing
\setlength{\parindent}{0pt}
\setlength{\parskip}{6pt}

\begin{document}

% Orange header
\begin{tcolorbox}[colback=osaorange,colframe=black,boxrule=2pt,arc=0pt,left=6pt,right=6pt,top=8pt,bottom=8pt]
\centering
\textcolor{white}{\textbf{\large ANSWER KEY: WEEKS 7, 8, 9}}\\
\textcolor{white}{\textbf{H524 Introduction to Biostatistics - Fall 2025}}\\
\textcolor{white}{\textbf{INSTRUCTOR USE ONLY}}
\end{tcolorbox}

\vspace{8pt}

\section*{Week 7: Nonparametric Methods (5 questions, 5 points)}

\subsection*{Question 1: When to Use Nonparametric Tests (1 point)}
\textbf{Correct Answer: \textcolor{correctcolor}{B}}\\
The data do not meet the assumptions required for parametric tests (e.g., normality)

\subsection*{Question 2: Sign Test vs. t-test (1 point)}
\textbf{Correct Answer: \textcolor{correctcolor}{B}}\\
Sign test\\
\textit{Explanation:} Pain score differences are highly skewed, violating normality assumption for paired t-test.

\subsection*{Question 3: Mann-Whitney U Test (1 point)}
\textbf{Correct Answer: \textcolor{correctcolor}{C}}\\
Independent samples t-test (two-sample t-test)

\subsection*{Question 4: Kruskal-Wallis Test (1 point)}
\textbf{Correct Answer: \textcolor{correctcolor}{B}}\\
Kruskal-Wallis test\\
\textit{Explanation:} Right-skewed data violates normality; Kruskal-Wallis is nonparametric alternative to one-way ANOVA.

\subsection*{Question 5: Ranking Data (1 point)}
\textbf{Correct Answer: \textcolor{correctcolor}{B}}\\
Tied values receive the average of the ranks they would have occupied

\vspace{10pt}

\noindent\rule{\textwidth}{0.4pt}

\section*{Week 8: Contingency Tables / Chi-Square (9 questions, 9 points)}

\subsection*{Question 1: Contingency Table Basics (1 point)}
\textbf{Correct Answer: \textcolor{correctcolor}{C}}\\
Two categorical variables

\subsection*{Question 2: Expected Frequencies (1 point)}
\textbf{Correct Answer: \textcolor{correctcolor}{A}}\\
(Row total × Column total) / Grand total

\subsection*{Question 3: Chi-Square Test Assumptions (1 point)}
\textbf{Correct Answer: \textcolor{correctcolor}{B}}\\
All expected cell frequencies should be at least 5

\subsection*{Question 4: Degrees of Freedom (1 point)}
\textbf{Correct Answer: \textcolor{correctcolor}{B}}\\
6\\
\textit{Calculation:} df = (r-1)(c-1) = (3-1)(4-1) = 2 × 3 = 6

\subsection*{Question 5: Interpreting Chi-Square Results (1 point)}
\textbf{Correct Answer: \textcolor{correctcolor}{B}}\\
There is a significant association between the variables\\
\textit{Explanation:} p = 0.032 $<$ 0.05, so we reject the null hypothesis of no association.

\subsection*{Question 6: Relative Risk (1 point)}
\textbf{Correct Answer: \textcolor{correctcolor}{B}}\\
The ratio of the risk of disease in exposed vs. unexposed groups

\subsection*{Question 7: Odds Ratio (1 point)}
\textbf{Correct Answer: \textcolor{correctcolor}{B}}\\
The odds of disease in the exposed group are 2.5 times the odds in the unexposed group

\subsection*{Question 8: 2×2 Table Setup (1 point)}
\textbf{Correct Answer: \textcolor{correctcolor}{C}}\\
Either variable can form rows\\
\textit{Explanation:} 2×2 tables are symmetric; either variable can be rows or columns.

\subsection*{Question 9: Fisher's Exact Test (1 point)}
\textbf{Correct Answer: \textcolor{correctcolor}{B}}\\
When expected cell frequencies are small (less than 5 in one or more cells)

\vspace{10pt}

\noindent\rule{\textwidth}{0.4pt}

\section*{Week 9: Correlation / Regression (7 questions, 7 points)}

\subsection*{Question 1: Correlation Coefficient Range (1 point)}
\textbf{Correct Answer: \textcolor{correctcolor}{C}}\\
-1 to 1

\subsection*{Question 2: Interpreting Correlation (1 point)}
\textbf{Correct Answer: \textcolor{correctcolor}{D}}\\
A strong negative linear relationship\\
\textit{Explanation:} $|r| = 0.85$ is strong; negative sign indicates negative relationship.

\subsection*{Question 3: Correlation vs. Causation (1 point)}
\textbf{Correct Answer: \textcolor{correctcolor}{B}}\\
The variables tend to increase together, but this does not prove causation

\subsection*{Question 4: Simple Linear Regression Equation (1 point)}
\textbf{Correct Answer: \textcolor{correctcolor}{B}}\\
The slope of the regression line

\subsection*{Question 5: Coefficient of Determination (1 point)}
\textbf{Correct Answer: \textcolor{correctcolor}{A}}\\
0.36\\
\textit{Calculation:} R² = r² = (0.60)² = 0.36

\subsection*{Question 6: Interpreting R² (1 point)}
\textbf{Correct Answer: \textcolor{correctcolor}{A}}\\
64\% of the variance in y is explained by x

\subsection*{Question 7: Regression Assumptions (1 point)}
\textbf{Correct Answer: \textcolor{correctcolor}{D}}\\
The independent variable (x) must be normally distributed\\
\textit{Explanation:} This is NOT an assumption. The residuals must be normal, but x does not.

\vspace{10pt}

\noindent\rule{\textwidth}{0.4pt}

\vspace{8pt}

\section*{Grading Summary}

\begin{center}
\begin{tabular}{lc}
\toprule
\textbf{Week} & \textbf{MC Points} \\
\midrule
Week 7 & 5 points \\
Week 8 & 9 points \\
Week 9 & 7 points \\
\bottomrule
\textbf{Total} & \textbf{21 points} \\
\end{tabular}
\end{center}

\vspace{8pt}

\textbf{Note:} Each week also includes a group project check-in component:
\begin{itemize}[leftmargin=20pt]
\item Week 7: Team Plan (5 points) - completion-based
\item Week 8: Progress Update (1 point) - completion-based
\item Week 9: Presentation Readiness (3 points) - completion-based
\end{itemize}

\end{document}
